% from ega2-1-fr.tex
\oldpage[II]{169}
et $\mathcal{L}=\mathcal{O}_Y$, donc $\mathcal{S}=(A[T])^{\sim}$ et
$\bigoplus_{n\in\mathbf{Z}}\mathcal{L}^{\otimes n}
=(A[T,T^{-1}])^{\sim}$. Or $A[T,T^{-1}]$ est l'anneau de fractions
$A[T]_T$ de $A[T]$, donc (\hyperref[II.8.4.4.3-fr]{8.4.4.3}) identifie le
premier membre au préschéma induit par $C=\mathbf{V}(\mathcal{L})$ sur
l'ouvert $D(T)$~; le complémentaire $V(T)$ de cet ouvert dans $C$ est
l'espace sous-jacent au sous-préschéma fermé de $C$ défini par l'idéal
$TA[T]$, c'est-à-dire la section nulle de $C$, donc $E=D(T)$.

L'isomorphisme (\hyperref[II.8.4.4.2-fr]{8.4.4.2}) sera d'autre part
conséquence de la dernière assertion, $\mathbf{V}(\mathcal{L}^{-1})$ étant le
complémentaire du lieu à l'infini de sa fermeture projective et $\widehat{E}$
le complémentaire de la section nulle de la fermeture projective de
$C=\mathbf{V}(\mathcal{L})$. Or, ces fermetures projectives sont
respectivement $\mathbf{P}(\mathcal{L}^{-1}\oplus\mathcal{O}_Y)$ et
$\mathbf{P}(\mathcal{L}\oplus\mathcal{O}_Y)$~; mais on peut écrire
$\mathcal{L}\oplus\mathcal{O}_Y
=\mathcal{L}\otimes(\mathcal{L}^{-1}\oplus\mathcal{O}_Y)$. L'existence de
l'isomorphisme canonique cherché résulte alors de
(\hyperref[II.4.1.4-fr]{4.1.4}), et tout revient à voir que cet isomorphisme
échange sections nulles et lieux à l'infini. On peut pour cela se borner au
cas où $Y=\operatorname{Spec}(A)$ est affine,
$L=Ac$, $L^{-1}=Ac'$, l'isomorphisme canonique
$L\otimes L^{-1}\to A$ faisant correspondre à $c\otimes c'$ l'élément $1$
de $A$. Alors $\mathbf{S}(L\oplus A)$ est produit tensoriel de $A[z]$ et de
$\bigoplus_{n\geq0}Ac^{\otimes n}$, $\mathbf{S}(L^{-1}\oplus A)$ produit
tensoriel de $A[z]$ et de $\bigoplus_{n\geq0}Ac'^{\otimes n}$, et
l'isomorphisme défini dans (\hyperref[II.4.1.4-fr]{4.1.4}) fait correspondre à
$z^h\otimes c'^{\otimes(n-h)}$ l'élément
$z^{n-h}\otimes c^{\otimes h}$. Or, dans
$\mathbf{P}(\mathcal{L}^{-1}\oplus\mathcal{O}_Y)$, le lieu à l'infini est
l'ensemble des points où s'annule la section $z$ et la section nulle est
l'ensemble des points où s'annule la section $c'$~; comme on a des définitions
analogues pour $\mathbf{P}(\mathcal{L}\oplus\mathcal{O}_Y)$, notre conclusion
résulte aussitôt de l'explicitation qui précède.
\end{proof}

\subsection{Comportements fonctoriels}
\label{subsection:II.8.5-fr}

\begin{env}[8.5.1]
\label{II.8.5.1-fr}
Soient $Y$, $Y'$ deux préschémas, $q:Y'\to Y$ un morphisme,
$\mathcal{S}$ (resp. $\mathcal{S}'$) une $\mathcal{O}_Y$-Algèbre
(resp. une $\mathcal{O}_{Y'}$-Algèbre) graduée quasi-cohérente à degrés
\emph{positifs}. Considérons un $q$-morphisme d'Algèbres graduées
\begin{equation}
\label{II.8.5.1.1-fr}
  \varphi:\mathcal{S}\to\mathcal{S}'.
\tag{8.5.1.1}
\end{equation}
On sait (\hyperref[II.1.5.6-fr]{1.5.6}) qu'il lui correspond canoniquement un
morphisme
\[
  \Phi=\operatorname{Spec}(\varphi):
  \operatorname{Spec}(\mathcal{S}')\to\operatorname{Spec}(\mathcal{S})
\]
tel que le diagramme
\begin{equation}
\label{II.8.5.1.2-fr}
  \xymatrix{
    C'\ar[r]^{\Phi}\ar[d] & C\ar[d]\\
    Y'\ar[r]_{q} & Y
  }
\tag{8.5.1.2}
\end{equation}
où on a posé $C=\operatorname{Spec}(\mathcal{S})$,
$C'=\operatorname{Spec}(\mathcal{S}')$, soit commutatif. 
\emph{Supposons en outre que
$\mathcal{S}_0=\mathcal{O}_Y$, $\mathcal{S}'_0=\mathcal{O}_{Y'}$}~; soient
$\varepsilon:Y\to C$ et $\varepsilon':Y'\to C'$ les immersions canoniques
(\hyperref[II.8.3.2-fr]{8.3.2})~; on a alors un diagramme commutatif
\begin{equation}
\label{II.8.5.1.3-fr}
  \xymatrix{
    Y'\ar[r]^{q}\ar[d]_{\varepsilon'} & Y\ar[d]^{\varepsilon}\\
    C'\ar[r]_{\Phi} & C
  }
\tag{8.5.1.3}
\end{equation}

